\documentclass{article}
\usepackage[T1]{fontenc}
\usepackage{lmodern}
\usepackage{graphicx}
\usepackage{amsmath,amssymb}
\usepackage[a4paper,margin=2cm]{geometry}
\usepackage[absolute,overlay]{textpos}
\setlength{\TPHorizModule}{1mm}
\setlength{\TPVertModule}{1mm}
\usepackage[indonesian]{babel}
\usepackage{hyperref}
\setlength{\emergencystretch}{2em}
% unit: Halaman 72

\begin{document}

% === Halaman 72 (llm) ===

\section*{Ukuran blok dan kunci}

\begin{itemize}
    \item Ukuran blok pesan ($n$) yang diproses selama enkripsi/dekripsi selalu tetap.
    
    Contoh: DES ($n = 64$), AES ($n = 128, 192, 256$)
    
    \item Makin besar ukuran blok mengarah ke efek \textit{diffusion} yang lebih besar.
    
    \item Ukuran kunci ($m$) bisa sama dengan ukuran blok atau berbeda dengan ukuran blok
    
    Contoh: DES ($m = 56$), AES ($m = 128, 192, 256$)
    
    \item Makin besar ukuran kunci mengarah ke efek \textit{confusion} yang lebih besar dan lebih tahan terhadap serangan \textit{brute force}.
\end{itemize}

\end{document}
